% ===== PRÉAMBULE CANONIQUE — à coller VERBATIM en tête de CHAQUE .tex =====
\documentclass[11pt,a4paper]{article}
\usepackage[utf8]{inputenc}\usepackage[T1]{fontenc}\usepackage{lmodern}
\usepackage[french]{babel}
\usepackage{amsmath,amssymb,mathtools}\usepackage{icomma}
\usepackage[version=4]{mhchem}          % \ce{...} : équations & formules
\usepackage{siunitx}\sisetup{locale=FR} % \qty{}{}, \si{}, \ang{}
\usepackage{chemfig}                    % \chemfig{...} : structures développées
\usepackage{xcolor}\usepackage{colortbl}
\usepackage{tikz}\usepackage{pgfplots}\pgfplotsset{compat=1.18}
\usepackage[most]{tcolorbox}\usepackage{enumitem}\usepackage{array,booktabs}
\usepackage[a4paper,margin=2.2cm]{geometry}
\usetikzlibrary{arrows.meta,calc,angles,quotes,positioning,patterns,backgrounds,shapes.geometric}
\definecolor{primary}{RGB}{222,49,99}\definecolor{primarydark}{RGB}{150,22,55}
\definecolor{defcol}{RGB}{0,114,178}\definecolor{methcol}{RGB}{52,92,148}
\definecolor{excol}{RGB}{0,135,90}\definecolor{attncol}{RGB}{200,40,55}
\tcbset{boxbase/.style={enhanced,breakable,boxrule=0.9pt,arc=2.5pt,
  left=3mm,right=3mm,top=2.3mm,bottom=2.3mm,fonttitle=\bfseries,coltitle=white}}
% --- boîtes TUTORIEL (noms EXACTS, ne pas renommer) ---
\newtcolorbox{cle}[1][]{boxbase,colback=defcol!6,colframe=defcol,
  title={\textsc{À retenir}\ifx\relax#1\relax\relax\else\textmd{ : \itshape #1}\fi}}
\newtcolorbox{methode}[1][]{boxbase,colback=methcol!7,colframe=methcol,sharp corners,
  title={\textsc{Méthode}\ifx\relax#1\relax\relax\else\textmd{ --- \itshape #1}\fi}}
\newtcolorbox{exemplebox}[1][]{boxbase,colback=excol!5,colframe=excol,
  title={\textsc{Exemple}\ifx\relax#1\relax\relax\else\textmd{ : \itshape #1}\fi}}
\newtcolorbox{piege}[1][]{boxbase,colback=attncol!6,colframe=attncol,
  title={\textsc{Piège}\ifx\relax#1\relax\relax\else\textmd{ : \itshape #1}\fi}}
\newtcolorbox{remarque}[1][]{boxbase,colback=black!4,colframe=black!45,
  title={\textsc{Remarque}\ifx\relax#1\relax\relax\else\textmd{ : \itshape #1}\fi}}
% --- boîtes EXERCICE / CORRIGÉ (noms EXACTS, auto-comptées) ---
\newtcolorbox[auto counter]{exo}[1][]{boxbase,colback=primary!4,colframe=primary,
  title={\textsc{Exercice}~\thetcbcounter\ifx\relax#1\relax\relax\else\textmd{ --- \itshape #1}\fi}}
\newtcolorbox[auto counter]{corrige}[1][]{boxbase,colback=excol!4,colframe=excol,
  title={\textsc{Corrigé}~\thetcbcounter\ifx\relax#1\relax\relax\else\textmd{ --- \itshape #1}\fi}}
\newcommand{\R}{\mathbb{R}}\newcommand{\N}{\mathbb{N}}\newcommand{\Z}{\mathbb{Z}}\newcommand{\ds}{\displaystyle}
% (un \usepackage{fancyhdr}+en-tête est possible ; il est ignoré côté web)
% =========================================================================
\begin{document}
\begin{exo}[octet et duet]
Un gaz noble a sa couche de valence saturée : $8$ électrons (octet), ou $2$ pour l'hélium (duet).
\begin{enumerate}[label=\alph*)]
  \item Combien d'électrons de valence possède un gaz noble autre que l'hélium ?
  \item Combien en possède l'hélium ? Comment nomme-t-on cette configuration à $2$ électrons ?
  \item Pourquoi les gaz nobles sont-ils chimiquement très peu réactifs ?
\end{enumerate}
\end{exo}

\begin{exo}[reconnaître une famille]
Associe le nombre d'électrons de valence à la famille : alcalins, alcalino-terreux, halogènes ou gaz
nobles.
\begin{enumerate}[label=\alph*)]
  \item $1$ électron de valence
  \item $2$ électrons de valence
  \item $7$ électrons de valence
  \item $8$ électrons de valence
\end{enumerate}
\end{exo}

\begin{exo}[quel ion se forme ?]
D'après la règle de l'octet, indique l'ion (avec sa charge) que forme chaque atome.
\begin{enumerate}[label=\alph*)]
  \item sodium $\ce{Na}$ ($1$ e$^-$ de valence, métal)
  \item magnésium $\ce{Mg}$ ($2$ e$^-$ de valence, métal)
  \item oxygène $\ce{O}$ ($6$ e$^-$ de valence)
  \item chlore $\ce{Cl}$ ($7$ e$^-$ de valence)
\end{enumerate}
\end{exo}

\begin{exo}[sens des tendances périodiques]
Complète avec le bon sens de variation.
\begin{enumerate}[label=\alph*)]
  \item L'électronégativité augmente vers la \dots\ (droite / gauche) et vers le \dots\ (haut / bas).
  \item Le rayon atomique augmente vers la \dots\ et vers le \dots.
  \item L'énergie d'ionisation varie dans le même sens que \dots\ (l'électronégativité / le rayon atomique).
\end{enumerate}
\end{exo}

\begin{exo}[comparer deux éléments]
Désigne le bon élément et justifie par sa position dans le tableau.
\begin{enumerate}[label=\alph*)]
  \item le plus électronégatif : $\ce{Na}$ ou $\ce{Cl}$ (même période) ?
  \item le plus grand rayon atomique : $\ce{Li}$ ou $\ce{K}$ (même colonne) ?
  \item la plus grande énergie d'ionisation : $\ce{Li}$ ou $\ce{F}$ (même période) ?
\end{enumerate}
\end{exo}

\end{document}
